% !TEX program = xelatex
%!TEX TS-program = xelatex
% Bu dosya sjuoz_word2latex.py tarafından otomatik üretildi.
\documentclass[9pt]{extarticle}
\usepackage{multicol}
\raggedcolumns  % sütun boyu doldurulmaz: başlık ve paragraf aralarında büyük boşluk açılmaz
\usepackage[top=0.5in, bottom=0.75in, right=0.75in, left=0.75in, headheight=0.25in, includehead]{geometry}
\usepackage{fontspec}
\usepackage{amsmath}
\usepackage{amssymb}
\usepackage{graphicx}
% Denklem sütuna sığmazsa: önce 2, sonra 3 satıra bölünmüş biçimi; o da sığmazsa küçültülür
\newsavebox{\sjuozeqbox}
\newlength{\sjuozeqmax}
\newcommand{\sjuozeqset}{\setlength{\sjuozeqmax}{\linewidth}\addtolength{\sjuozeqmax}{-2.5em}}
\newcommand{\sjuozeqput}{\ifdim\wd\sjuozeqbox>\sjuozeqmax
  \typeout{SJUOZ-EQ-SCALED: \theequation; \the\sjuozeqmax; \the\wd\sjuozeqbox}% kontrol raporu için
  \resizebox{\sjuozeqmax}{!}{\usebox{\sjuozeqbox}}\else\usebox{\sjuozeqbox}\fi}
\newcommand{\sjuozfit}[3]{\sjuozeqset\sbox{\sjuozeqbox}{$\displaystyle #1$}%
  \ifdim\wd\sjuozeqbox>\sjuozeqmax\sbox{\sjuozeqbox}{$\displaystyle #2$}\fi
  \ifdim\wd\sjuozeqbox>\sjuozeqmax\sbox{\sjuozeqbox}{$\displaystyle #3$}\fi\sjuozeqput}
\newcommand{\sjuozfitone}[1]{\sjuozeqset\sbox{\sjuozeqbox}{$\displaystyle #1$}\sjuozeqput}
\usepackage{xcolor}
\usepackage{array}
\usepackage{tabularx}
\usepackage{xltabular}
\usepackage{pdflscape}
\usepackage{booktabs}
\usepackage{multirow}
\usepackage{caption}
\usepackage{subcaption}
\usepackage{titlesec}
\usepackage{fancyhdr}
\usepackage{qrcode}
\usepackage{xurl}
\usepackage{lastpage}
\usepackage{indentfirst}
\usepackage{eso-pic}
\usepackage{tikz}
\usepackage[colorlinks=true, urlcolor=blue, linkcolor=black, citecolor=black]{hyperref}

\titleformat{\section}{\large\bfseries\centering}{\thesection.}{1.0em}{}
\titleformat{\subsection}{\large\bfseries}{\thesubsection.}{0.5em}{}
\titleformat{\subsubsection}{\large\bfseries}{\thesubsubsection.}{0.5em}{}
\titlespacing*{\section}{0pt}{8pt}{6pt}
\titlespacing*{\subsection}{0pt}{8pt}{4pt}
\titlespacing*{\subsubsection}{0pt}{6pt}{3pt}

\newcommand{\titlesize}{\fontsize{12}{10}}
\renewcommand{\normalsize}{\fontsize{9pt}{9pt}\selectfont}
\renewcommand{\large}{\fontsize{10pt}{10pt}\selectfont}
\renewcommand{\small}{\fontsize{8pt}{8pt}\selectfont}
\DeclareCaptionLabelFormat{sjuozparen}{#1 (#2)}
\captionsetup[table]{labelfont=bf, justification=raggedright, singlelinecheck=false}
\captionsetup[figure]{labelfont=bf, justification=centering, singlelinecheck=false}
\captionsetup[sub]{labelformat=parens, labelfont=bf, font=small, justification=centering}
\definecolor{sjuozblue}{HTML}{C1D1EB}
\definecolor{colorabs}{HTML}{E2EAF4}

% Yazı tipleri: Times New Roman / Garamond yoksa benzerleri kullanılır
\IfFontExistsTF{Times New Roman}{\setmainfont{Times New Roman}}{%
  \IfFontExistsTF{TeX Gyre Termes}{\setmainfont{TeX Gyre Termes}}{%
    \IfFontExistsTF{Liberation Serif}{\setmainfont{Liberation Serif}}{}}}
\IfFontExistsTF{Garamond}{\newfontfamily{\garamond}{Garamond}}{%
  \IfFontExistsTF{EB Garamond}{\newfontfamily{\garamond}{EB Garamond}}{\newcommand{\garamond}{\rmfamily}}}

\newcommand{\sjuozdoi}{https://doi.org/10.25271/sjuoz.2026.14.4.2006}
% İlk sayfadaki bilgi kutusunun sayfa altından yüksekliği ve metin için ayrılan boşluk
\newcommand{\firstfooterbottom}{0.4in}
\newcommand{\firstpagereserve}{1.4in}
% Sadece ilk sayfada metin alanını kısalt; ilk sayfa basıldıktan sonra eski haline döndür
\makeatletter
\newlength{\sjuozfulltextheight}
\AtBeginDocument{%
  \setlength{\sjuozfulltextheight}{\textheight}%
  \global\advance\textheight by -\firstpagereserve
  \global\@colht\textheight \global\@colroom\textheight \global\vsize\textheight
  \AddToHookNext{shipout/after}{\global\textheight\sjuozfulltextheight}}
\makeatother

% Başlık ve altlık (header/footer)
\pagestyle{fancy}
\fancyhf{}
\renewcommand{\headrulewidth}{0pt}
\lhead{\large \textit{Taha \& Saleem}}
\rhead{\small{\garamond{\textbf{SJUOZ|VOL14|OCT 2026|Pxxx-xxx}}}}
\lfoot{\raisebox{8pt}{\textit{DOI:}\href{\sjuozdoi}{\textit{\small{\sjuozdoi}}}}}
\cfoot{\raisebox{8pt}{\thepage}}

\fancypagestyle{firstpage}{%
  \fancyhf{}
  \renewcommand{\headrulewidth}{0pt}
  \rhead{\IfFileExists{SJUOZ.png}{\includegraphics[width=1.04in]{SJUOZ.png}}{\textbf{SJUOZ}}}
  \lhead{\colorbox{sjuozblue}{%
      \begin{minipage}[b]{0.83\textwidth}
        \vspace{0.15cm}
        \hspace{-3pt}\ \normalsize \noindent \garamond{Science Journal of University of Zakho}\\
        \normalsize \textit{\garamond{DOI: \sjuozdoi}}
        \vspace{0.2cm}
      \end{minipage}}}
}

% İlk sayfanın altındaki bilgi kutusu (sayfaya mutlak konumla yerleştirilir)
\newcommand{\sjuozfirstfooter}{%
  \begin{minipage}[b]{\textwidth}
      \rule{\textwidth}{1.5pt} \\[0.15cm]
      \begin{minipage}[t]{0.2\textwidth}\hspace{1pt} \small Access this article online \\
        \vspace{0.1cm}
        {\hypersetup{urlcolor=black}\qrcode[height=2.6162cm]{https://sjuoz.uoz.edu.krd/index.php/sjuoz/article/view/2006}}
      \end{minipage}\begin{minipage}[t]{0.4\textwidth}
        \small DOI:\href{\sjuozdoi}{\sjuozdoi} \\ \\[0.5cm]
        Printed ISSN 2663-628X; \\
        Electronic ISSN 2663-6298
      \end{minipage}\begin{minipage}[t]{0.4\textwidth}\small \textbf{Science Journal of University of Zakho} \\
        \textbf{Vol. 14, No. 04, pp. xxx-xxx, October-2026} \\[0.5cm]
        This is an open access article under a CC BY-NC-SA 4.0 license \\
        (\href{https://creativecommons.org/licenses/by-nc-sa/4.0/}{https://creativecommons.org/licenses/by-nc-sa/4.0/})
      \end{minipage} \\[0.15cm]
      \rule{\textwidth}{1.5pt} \\[0.15cm]
      \centering \normalsize \thepage
  \end{minipage}}

% ORCID simgesi: orcid.png varsa o, yoksa çizilmiş yeşil "iD" simgesi
\definecolor{orcidgreen}{HTML}{A6CE39}
\newcommand{\orcidicon}{\IfFileExists{orcid.png}{\includegraphics[width=9pt]{orcid.png}}{%
  \tikz[baseline=-0.7ex]{\fill[orcidgreen] (0,0) circle (4.5pt);
  \node[text=white, font=\fontsize{4.5pt}{4.5pt}\selectfont\bfseries, inner sep=0pt] at (0,0) {iD};}}}
\newcommand{\orcidlink}[1]{\href{#1}{\orcidicon}}

\begin{document}
\setcounter{page}{1}
\thispagestyle{firstpage}
\AddToShipoutPictureFG*{\put(\LenToUnit{0.75in},\LenToUnit{\firstfooterbottom}){\sjuozfirstfooter}}

\noindent \large \textbf{Original Article} \\[-5pt]
\begin{center} \titlesize \textbf{WIDE TABLES ON LANDSCAPE PAGES} \end{center}
\begin{center} \large \textbf{Hevi Taha \textsuperscript{1,*} and Azad Saleem \textsuperscript{1}}
\end{center}
\noindent
\large \textsuperscript{1}Department of Physics, University of Zakho, Duhok, Iraq.\\

\begin{center} \textsuperscript{*}\textit{Corresponding author, E-mail:}\href{mailto:hevi.taha@uoz.edu.krd}{hevi.taha@uoz.edu.krd}
\end{center}
\rule{\textwidth}{1.5pt}
\begin{minipage}[t]{0.175\textwidth}
  \hspace{0cm} \vspace{0cm}
\end{minipage} \begin{minipage}[t]{0.8\textwidth}
\large \textbf{ABSTRACT}\end{minipage} \\[-6pt]
\rule{\textwidth}{1.5pt} \\
\begin{minipage}[t]{0.175\textwidth}
  \vspace{1cm} Received \\ xxx \\[0.5cm] Accepted \\ xxx \\[0.5cm] Published \\ xxx \\[0.25cm]
\end{minipage} \colorbox{colorabs}{\begin{minipage}[t]{0.81\textwidth}
  \large{Wide tables are placed in a landscape section in Word, or are simply too wide for a portrait page; both must be printed on rotated pages without squeezing the columns.
  \\[2pt]
  \normalsize \textbf{Keywords:} landscape; tables.
  }
\end{minipage}}
\\[-10pt]

\normalsize
\noindent \rule{\textwidth}{0.75pt}
\begin{multicols}{2}
\section{INTRODUCTION}
\hyperlink{tab:1}{Table 1} is in a landscape section of the Word file and \hyperlink{fig:1}{Figure 1} is next to it. \hyperlink{tab:2}{Table 2} has twelve columns in a normal portrait section (\hyperlink{ref:1}{Taha, 2025}).

\end{multicols}
\begin{landscape}
\noindent
\begin{minipage}{\linewidth}
\centering
\hypertarget{tab:1}{}
\captionof{table}{Measurements in a landscape section.}
\vspace{-6pt}
\footnotesize
\begin{tabularx}{\linewidth}{@{}>{\raggedright\arraybackslash}X>{\centering\arraybackslash}X>{\centering\arraybackslash}X>{\centering\arraybackslash}X>{\centering\arraybackslash}X>{\centering\arraybackslash}X@{}}
\toprule
\textbf{Sample} & \textbf{Temperature} & \textbf{Pressure} & \textbf{Humidity} & \textbf{Wind speed} & \textbf{Notes about the measurement} \\ \midrule
A1 & 21.5 & 1012 & 40 & 3.2 & Morning, clear sky, calibrated sensor \\ \addlinespace[0.3em]
A2 & 25.1 & 1009 & 35 & 4.8 & Noon, light clouds, calibrated sensor \\ \addlinespace[0.3em]
A3 & 19.8 & 1015 & 55 & 2.1 & Evening, humid, sensor recalibrated \\
\bottomrule
\end{tabularx}
\end{minipage}
\par\vspace{6pt}

\noindent\begin{minipage}{\linewidth}
\centering\hypertarget{fig:1}{}
\includegraphics[width=\linewidth, height=0.45\textheight, keepaspectratio]{figures/Figure1.IMG}
\captionof{figure}{A picture in the landscape section.}
\end{minipage}
\par\vspace{6pt}

\end{landscape}
\begin{multicols}{2}
The text continues on portrait pages after the landscape section.

\end{multicols}
\begin{landscape}
\noindent
\begin{minipage}{\linewidth}
\centering
\hypertarget{tab:2}{}
\captionof{table}{Twelve features.}
\vspace{-6pt}
\footnotesize
\begin{tabularx}{\linewidth}{@{}>{\raggedright\arraybackslash}X>{\centering\arraybackslash}X>{\centering\arraybackslash}X>{\centering\arraybackslash}X>{\centering\arraybackslash}X>{\centering\arraybackslash}X>{\centering\arraybackslash}X>{\centering\arraybackslash}X>{\centering\arraybackslash}X>{\centering\arraybackslash}X>{\centering\arraybackslash}X>{\centering\arraybackslash}X@{}}
\toprule
\textbf{Feature 1} & \textbf{Feature 2} & \textbf{Feature 3} & \textbf{Feature 4} & \textbf{Feature 5} & \textbf{Feature 6} & \textbf{Feature 7} & \textbf{Feature 8} & \textbf{Feature 9} & \textbf{Feature 10} & \textbf{Feature 11} & \textbf{Feature 12} \\ \midrule
0.78 & 0.89 & 1.00 & 1.11 & 1.22 & 1.33 & 1.44 & 1.56 & 1.67 & 1.78 & 1.89 & 2.00 \\ \addlinespace[0.3em]
1.56 & 1.67 & 1.78 & 1.89 & 2.00 & 2.11 & 2.22 & 2.33 & 2.44 & 2.56 & 2.67 & 2.78 \\ \addlinespace[0.3em]
2.33 & 2.44 & 2.56 & 2.67 & 2.78 & 2.89 & 3.00 & 3.11 & 3.22 & 3.33 & 3.44 & 3.56 \\ \addlinespace[0.3em]
3.11 & 3.22 & 3.33 & 3.44 & 3.56 & 3.67 & 3.78 & 3.89 & 4.00 & 4.11 & 4.22 & 4.33 \\ \addlinespace[0.3em]
3.89 & 4.00 & 4.11 & 4.22 & 4.33 & 4.44 & 4.56 & 4.67 & 4.78 & 4.89 & 5.00 & 5.11 \\ \addlinespace[0.3em]
4.67 & 4.78 & 4.89 & 5.00 & 5.11 & 5.22 & 5.33 & 5.44 & 5.56 & 5.67 & 5.78 & 5.89 \\ \addlinespace[0.3em]
5.44 & 5.56 & 5.67 & 5.78 & 5.89 & 6.00 & 6.11 & 6.22 & 6.33 & 6.44 & 6.56 & 6.67 \\ \addlinespace[0.3em]
6.22 & 6.33 & 6.44 & 6.56 & 6.67 & 6.78 & 6.89 & 7.00 & 7.11 & 7.22 & 7.33 & 7.44 \\
\bottomrule
\end{tabularx}
\end{minipage}
\par\vspace{6pt}

\end{landscape}
\begin{multicols}{2}
\section{CONCLUSIONS}
Both tables must be readable (\hyperlink{ref:1}{Taha, 2025}).

\section*{REFERENCES}
\begin{list}{}{
    \setlength{\leftmargin}{1.5em}
    \setlength{\itemindent}{-1.5em}
    \setlength{\parsep}{0pt}
    \setlength{\itemsep}{6pt}
}
\item \hypertarget{ref:1}{}Taha, H. (2025). Rotated pages in journals. \textit{Journal of Layout}, 4(2), 10–12.

\end{list}
\end{multicols}

\end{document}
